\begin{theorem}[7.6.4]
\EGAReaderTarget{ega4.c04.thm.7-6-4}{IV.2.s7.6.4-fr}
(Zariski-Nagata). --- Sea $A$ un anillo semilocal noetheriano. Las
condiciones siguientes son equivalentes:
\begin{enumerate}
\item[\emph{a)}] Para toda $A$-álgebra finita reducida $C$, el completado
$\widehat C$ es un anillo reducido.
\item[\emph{a$'$)}] Para todo anillo cociente íntegro $B$ de $A$, de cuerpo
de fracciones $K$, el completado $\widehat B$ es reducido, y los cuerpos
componentes $L_i$ del anillo total de fracciones de $\widehat B$ son
extensiones separables de $K$.
\item[\emph{a$''$)}] Las fibras formales de $A$ son geométricamente
reducidas \textup{(es decir, para todo anillo cociente íntegro $B$ de $A$, de
cuerpo de fracciones $K$, la $K$-álgebra $\widehat B\otimes_BK$ es
separable (4.6.2))}.
\item[\emph{b)}] Todo anillo cociente íntegro $B$ de $A$ es un anillo
japonés.
\end{enumerate}
\end{theorem}

\begin{proof}
Para demostrar que \emph{a)} implica \emph{a$'$)}, basta comprobar que los $L_i$
son extensiones separables de $K$, o también (4.6.1) que para toda
extensión finita $K'$ de $K$, el anillo $\widehat B\otimes_BK'$ es
reducido. Ahora bien, $K'$ está generado por un número finito de
elementos enteros sobre $B$, y estos elementos generan una
sub-$B$-álgebra finita $B'$ de $K'$, cuyo cuerpo de fracciones es
$K'$. Se tiene $\widehat{B}'=\widehat B\otimes_BB'$
((\textbf 0$_{\mathrm I}$, 7.3.3) y Bourbaki, \emph{Alg. comm.}, cap. IV, § 2,
n\textsuperscript{o} 5, cor. 3 de
\oldpage[IV-2]{210}
la prop. 9), por lo que
$\widehat B\otimes_BK'=\widehat{B}'\otimes_{B'}K'$ por la
asociatividad del producto tensorial; pero $B'$ es una $A$-álgebra
íntegra finita, de modo que $\widehat{B}'$ es un anillo reducido
en virtud de $a)$. Como $\widehat{B}'$ es un $B'$-módulo plano,
los elementos $\neq0$ de $B'$ son $\widehat{B}'$-regulares, lo que
implica que $\widehat{B}'\otimes_{B'}K'$ se identifica con un
subanillo del anillo total de fracciones de $\widehat{B}'$ y
\emph{a fortiori} es reducido.

Para ver que $a')$ implica $a'')$, consideremos un punto cualquiera
$x$ de $X=\operatorname{Spec}(A)$, y sea $Y$ el subesquema cerrado
íntegro de $X$ cuyo espacio subyacente es $\overline{\{x\}}$;
se tiene $Y=\operatorname{Spec}(B)$, donde $B$ es un anillo
cociente íntegro de $A$, y
$\operatorname{Spec}(\widehat B)=Y\times_X\operatorname{Spec}(\widehat A)$,
por lo que la fibra formal de $A$ en el punto $x$ es la misma que
la de $B$, o también es igual a
$\operatorname{Spec}(\widehat B\otimes_BK)$. Como los anillos
locales de $\operatorname{Spec}(\widehat B\otimes_BK)$ son los
de $\operatorname{Spec}(\widehat B)$ en los puntos de la fibra
de $x$, la hipótesis implica que $\widehat B\otimes_BK$ es
reducido, y las condiciones de $a')$ implican, por tanto, que
$\widehat B\otimes_BK$ es una $K$-álgebra separable (4.6.1).

La condición $a'')$ implica $a)$, pues de (7.3.15) se deduce que
las fibras formales de $C$ son entonces geométricamente reducidas,
y, si $C$ es reducido, también lo es $\widehat C$ (caso
particular de (7.3.17)).

El hecho de que $a')$ implique $b)$ es un caso particular de
(\textbf 0, 23.1.7). Queda, por tanto, por demostrar que $b)$ implica
$a)$. Observemos que, si $C$ es una $A$-álgebra finita, para
todo ideal primo $\mathfrak q$ de $C$, la imagen inversa
$\mathfrak p$ de $\mathfrak q$ en $A$ es un ideal primo, y
$C/\mathfrak q$ es una $(A/\mathfrak p)$-álgebra finita; por
tanto, la hipótesis $b)$ implica que todo anillo cociente
íntegro de $C$ es un anillo japonés. Se reduce así la cuestión
a demostrar que, bajo la hipótesis $b)$, \emph{el completado de todo
cociente íntegro de $A$ es reducido}. Razonemos por inducción
sobre $n=\dim(A)$; la afirmación es trivial para $n=0$.
Sustituyendo $A$ por los cocientes de $A$ por sus ideales primos
minimales $\mathfrak p_i$, es posible limitarse al caso en que
$A$ es íntegro (todo cociente íntegro de $A$ es cociente de
uno de los $A/\mathfrak p_i$). Para todo ideal primo
$\mathfrak p\neq0$, la hipótesis de inducción muestra ya que
el completado de $A/\mathfrak p$ es reducido; basta, por
tanto, demostrar que $\widehat A$ es reducido. Además,
el cierre entero $A'$ de $A$ es, por hipótesis, un
$A$-módulo finitamente generado y, por tanto, un anillo
semilocal noetheriano; $\widehat A$ se identifica con un
subanillo de $\widehat{A}'$ ((\textbf 0$_{\mathrm I}$, 7.3.3) y
Bourbaki, \emph{Alg. comm.}, cap. IV, § 2, n\textsuperscript{o} 5,
cor. 3 de la prop. 9). Bastará, pues, demostrar que
$\widehat{A}'$ es reducido. Se ha visto más arriba que la hipótesis
$b)$ también se verifica para $A'$, que tiene, además, dimensión
$n$ (\textbf 0, 16.1.5); es posible, por tanto, limitarse al caso en
que $A$ es \emph{integralmente cerrado}. Sea $t\neq0$ un elemento
del radical de Jacobson de $A$, y sean $\mathfrak q_j$
($1\leq j\leq n$) los ideales primos minimales entre los
que contienen a $tA$; se tienen las propiedades siguientes:
\begin{enumerate}
\item[(i)] $t$ es regular.
\item[(ii)] $A/tA$ no tiene ideales primos asociados inmersos.
\item[(iii)] Los $A_{\mathfrak q_j}$ son anillos de valoración
discreta.
\item[(iv)] Los completados de los $A/\mathfrak q_j$ son
reducidos.
\end{enumerate}
En efecto, (i) es trivial porque $A$ es íntegro y $t\neq0$.
Como $A$ es integralmente cerrado, $A/tA$ satisface
$(\mathrm S_1)$, es decir (5.7.7), no tiene ideales primos
asociados inmersos (Bourbaki, \emph{Alg. comm.}, cap. VII,
§ 1, n\textsuperscript{o} 4, prop. 8). También por ser $A$ integralmente
\oldpage[IV-2]{211}
cerrado, lo son los $A_{\mathfrak q_j}$, y se
sabe (\emph{loc. cit.}) que estos anillos tienen dimensión 1;
son, por tanto, anillos de valoración discreta, de donde se
deduce (iii). Finalmente, (iv) se sigue de la hipótesis de
inducción. La demostración quedará concluida si probamos
el siguiente

\begin{lemma}[7.6.4.1]
\label{IV.2.s7.6.4.1-fr}
(Zariski). --- Sean $A$ un anillo semilocal noetheriano y $t$
un elemento de su radical de Jacobson que satisface las
condiciones \textup{(i)} a \textup{(iv)} anteriores; entonces $\widehat A$
es reducido.
\end{lemma}
